\documentclass[11pt,a4paper]{article}
\usepackage[margin=1in]{geometry}
\usepackage{amsmath,amssymb,amsfonts}
\usepackage{times}

\title{The Contexts-Systems-Modalities (CSM) framework of quantum logic}
\author{Karl Svozil\\
TU Wien}
\date{}

\begin{document}

\maketitle

\begin{abstract}

The Contexts-Systems-Modalities (CSM) framework, developed by Auff\`eves and Grangier, introduces an ontology for quantum mechanics built on three primitive notions: contexts (macroscopic measurement arrangements), systems (the entities under study), and modalities (sets of definite, repeatable outcomes associated with a given system--context pair). Within this approach, a modality represents an objective property that can be predicted with certainty inside its specific context. Superpositions are reinterpreted as modalities belonging to mutually incompatible contexts.

The CSM framework shares significant structural similarities with traditional quantum logic \`a la Birkhoff and von Neumann, particularly in its departure from classical Boolean logic, as well as with the operational empirical logic of Foulis and Randall, which emphasises collections of compatible tests and explicit contextuality. Like Foulis--Randall's manual-based approach, CSM places strong emphasis on contextual outcomes. However, it seeks to provide an explicit realist interpretation in which modalities are regarded as contextually objective properties.

A central feature of CSM is a quantization axiom limiting the number of mutually exclusive modalities within any single context. Combined with an assumption of extra-contextuality (relating modalities across different contexts), the quantum formalism---including vector-space (Hilbert space) structure and the Born rule---is recovered by consistently assuming vector space entities for the modalities.

Compared to Birkhoff--von Neumann quantum logic, which works directly with the orthomodular lattice of projections, and to the more neutral operational stance of Foulis and Randall, the CSM approach attempts to bridge operational structures with a contextual form of realism.

\end{abstract}

\vspace{0.5cm}
\textbf{Keywords:} quantum logic, CSM framework, modalities, contextuality, Birkhoff--von Neumann, Foulis--Randall, operational quantum logic

\end{document}
